%%%%%%%%%%%%%%%%%%%%%%%%%%%%%%%%%%%%%%%%%%%%%%%%%%%%%%%%%%%%%%%%%%%%%%%%%%%%%%%%%%
\begin{frame}[fragile]\frametitle{}
\begin{center}
{\Large Implementation}
\end{center}
\end{frame}

%%%%%%%%%%%%%%%%%%%%%%%%%%%%%%%%%%%%%%%%%%%%%%%%%%%%%%%%%%%%%%%%%%%%%%%%%%%%%%%%%%
\begin{frame}[fragile]\frametitle{}
\begin{center}
{\Large Node Implementation Examples}
\end{center}
\end{frame}

%%%%%%%%%%%%%%%%%%%%%%%%%%%%%%%%%%%%%%%%%%%%%%%%%%%%%%%%%%%
\begin{frame}[fragile]\frametitle{Example: Node with RAG}
% [code-font: forced \tiny | overflow at the default size]
      \begin{lstlisting}[language=Python, basicstyle=\tiny\ttfamily]
from langchain_community.vectorstores import FAISS
from langchain_groq import ChatGroq
from langgraph.graph import MessagesState

class RAGState(MessagesState):
    retrieved_docs: list

# Initialize RAG components (embeddings: any LangChain embeddings object)
vectorstore = FAISS.load_local("docs_index", embeddings,
                               allow_dangerous_deserialization=True)
llm = ChatGroq(model="openai/gpt-oss-20b")

def rag_node(state: RAGState) -> dict:
    """Node that performs RAG operation"""
    # 1. Retrieve relevant documents
    query = state["messages"][-1].content
    docs = vectorstore.similarity_search(query, k=3)
    context = "\n".join(doc.page_content for doc in docs)

    # 2. Generate answer with context
    prompt = f"""Context: {context}

Question: {query}

Answer based on the context provided:"""
    response = llm.invoke(prompt)

    # 3. Return ONLY what changed; add_messages appends the answer
    return {"messages": [response], "retrieved_docs": docs}
      \end{lstlisting}
\end{frame}

%%%%%%%%%%%%%%%%%%%%%%%%%%%%%%%%%%%%%%%%%%%%%%%%%%%%%%%%%%%
\begin{frame}[fragile]\frametitle{Example: Node with API Call}
 Any REST API, database query, or external service can be called from nodes

% [code-font: forced \scriptsize | overflow at the default size]
      \begin{lstlisting}[language=Python, basicstyle=\scriptsize\ttfamily]
import os
import requests
from pydantic import BaseModel

class WeatherState(BaseModel):
    location: str
    weather: str = ""

def weather_api_node(state: WeatherState) -> dict:
    """Node that calls external weather API"""
    # Keys come from the environment, never from source code
    api_key = os.getenv("OPENWEATHER_API_KEY")
    response = requests.get(
        "https://api.openweathermap.org/data/2.5/weather",
        params={"q": state.location, "appid": api_key, "units": "metric"},
        timeout=20,
    )
    response.raise_for_status()
    data = response.json()

    # Process and return only the changed field
    weather_info = (f"Temperature: {data['main']['temp']} deg C, "
                    f"Conditions: {data['weather'][0]['description']}")
    return {"weather": weather_info}
      \end{lstlisting}

\end{frame}

%%%%%%%%%%%%%%%%%%%%%%%%%%%%%%%%%%%%%%%%%%%%%%%%%%%%%%%%%%%
\begin{frame}[fragile]\frametitle{Example: Node as Tool-Calling Agent}

      \begin{itemize}
        \item Node delegates to inner agent for complex tool-using tasks
        \item Outer graph manages high-level workflow
      \end{itemize}

% [code-font: forced \tiny | overflow at the default size]
      \begin{lstlisting}[language=Python, basicstyle=\tiny\ttfamily]
from langchain.agents import create_agent
from langchain_core.messages import AIMessage
from langchain_core.tools import tool
from pydantic import BaseModel

@tool
def search(query: str) -> str:
    """Search the web."""
    return tavily_search(query)

@tool
def calculator(expression: str) -> str:
    """Perform calculations."""
    return str(safe_eval(expression))

research_agent = create_agent(model=llm, tools=[search, calculator])

class ResearchState(BaseModel):
    current_task: str
    research_results: str = ""

def research_node(state: ResearchState) -> dict:
    """Node that is itself an agent with tools"""
    # This node runs a full agent loop internally
    result = research_agent.invoke({"messages": [("user", state.current_task)]})
    answer = result["messages"][-1].content
    return {"research_results": answer}
      \end{lstlisting}

\end{frame}


%%%%%%%%%%%%%%%%%%%%%%%%%%%%%%%%%%%%%%%%%%%%%%%%%%%%%%%%%%%%%%%%%%%%%%%%%%%%%%%%%%
\begin{frame}[fragile]\frametitle{}
\begin{center}
{\Large End-to-End example}
\end{center}
\end{frame}

%%%%%%%%%%%%%%%%%%%%%%%%%%%%%%%%%%%%%%%%%%%%%%%%%%%%%%%%%%%
\begin{frame}[fragile]\frametitle{Implementing a Complete Tool-Calling Agent}
      \begin{itemize}
        \item Understand user input (think)
        \item Select the appropriate tool
        \item Execute tool calls
        \item Generate the final response
        \item The state object carries every decision from node to node
      \end{itemize}
\end{frame}


%%%%%%%%%%%%%%%%%%%%%%%%%%%%%%%%%%%%%%%%%%%%%%%%%%%%%%%%%%%
\begin{frame}[fragile]\frametitle{Define State}
      \begin{lstlisting}[language=Python]
from typing import Annotated, Literal, Optional
from langchain_core.messages import AnyMessage
from langgraph.graph.message import add_messages
from pydantic import BaseModel, Field

class AgentState(BaseModel):
    """The context object that flows through every node."""
    messages: Annotated[list[AnyMessage], add_messages] = Field(default_factory=list)
    current_input: str
    thought: str = ""
    selected_tool: Optional[str] = None
    tool_input: str = ""
    tool_output: str = ""
    final_answer: str = ""
    status: Literal["STARTING", "NEED_TOOL", "GENERATE_RESPONSE",
                    "SUCCESS", "ERROR"] = "STARTING"
      \end{lstlisting}

\end{frame}

%%%%%%%%%%%%%%%%%%%%%%%%%%%%%%%%%%%%%%%%%%%%%%%%%%%%%%%%%%%
\begin{frame}[fragile]\frametitle{Define Tools}
      \begin{lstlisting}[language=Python]
from langchain_core.tools import tool
from langchain_community.tools import WikipediaQueryRun
from langchain_community.utilities import WikipediaAPIWrapper

@tool
def calculator(expression: str) -> str:
    """Used for performing mathematical calculations, e.g. '25 * 4 + 10'."""
    allowed = set("0123456789+-*/(). ")
    if not set(expression) <= allowed:
        return "Invalid expression"
    return str(eval(expression, {"__builtins__": {}}))

wikipedia = WikipediaQueryRun(api_wrapper=WikipediaAPIWrapper())  # name: "wikipedia"

# Available tools, looked up by the name the model chooses
tools = {t.name: t for t in (calculator, wikipedia)}
      \end{lstlisting}

\end{frame}


%%%%%%%%%%%%%%%%%%%%%%%%%%%%%%%%%%%%%%%%%%%%%%%%%%%%%%%%%%%
\begin{frame}[fragile]\frametitle{Implement Core Nodes: Think}
\begin{lstlisting}[language=Python]
from langchain_groq import ChatGroq

llm = ChatGroq(model="openai/gpt-oss-20b", temperature=0)

class Decision(BaseModel):
    thought: str = Field(description="Reasoning about the next action")
    need_tool: bool
    tool: Optional[str] = Field(None, description="Tool name if needed")
    tool_input: str = ""

def think(state: AgentState) -> dict:
    """Think about the next action"""
    tool_list = [f"{t.name}: {t.description}" for t in tools.values()]
    prompt = (f"User input: {state.current_input}\n"
              f"Available tools: {tool_list}\n"
              "Decide whether a tool is needed, which one, and with what input.")
    d = llm.with_structured_output(Decision).invoke(prompt)   # validated JSON
    return {"thought": d.thought,
            "selected_tool": d.tool,
            "tool_input": d.tool_input,
            "status": "NEED_TOOL" if d.need_tool else "GENERATE_RESPONSE"}
\end{lstlisting}

\end{frame}

%%%%%%%%%%%%%%%%%%%%%%%%%%%%%%%%%%%%%%%%%%%%%%%%%%%%%%%%%%%
\begin{frame}[fragile]\frametitle{Implement Core Nodes: Execute Tool}
\begin{lstlisting}[language=Python]
def execute_tool(state: AgentState) -> dict:
    """Execute tool call"""
    tool_ = tools.get(state.selected_tool)
    if tool_ is None:
        return {"status": "ERROR", "thought": "Selected tool not found"}
    try:
        result = tool_.invoke(state.tool_input)
        return {"tool_output": str(result), "status": "GENERATE_RESPONSE"}
    except Exception as e:
        return {"status": "ERROR", "thought": f"Tool execution failed: {e}"}
\end{lstlisting}

\end{frame}


%%%%%%%%%%%%%%%%%%%%%%%%%%%%%%%%%%%%%%%%%%%%%%%%%%%%%%%%%%%
\begin{frame}[fragile]\frametitle{Implement Core Nodes: Generate Response}
\begin{lstlisting}[language=Python]
from langchain_core.messages import AIMessage

def generate_response(state: AgentState) -> dict:
    """Generate the final response"""
    prompt = f"""Generate a clear and helpful response for the user.
User input: {state.current_input}
Thought process: {state.thought}
Tool output: {state.tool_output}"""
    response = llm.invoke(prompt)
    return {"final_answer": response.content,
            "messages": [AIMessage(content=response.content)],
            "status": "SUCCESS"}
\end{lstlisting}

\end{frame}

%%%%%%%%%%%%%%%%%%%%%%%%%%%%%%%%%%%%%%%%%%%%%%%%%%%%%%%%%%%
\begin{frame}[fragile]\frametitle{Build the Complete Workflow}
% [code-font: forced \scriptsize | overflow at the default size]
\begin{lstlisting}[language=Python, basicstyle=\scriptsize\ttfamily]
from langgraph.graph import StateGraph, START, END

def route_after_think(state: AgentState) -> str:
    return "execute_tool" if state.status == "NEED_TOOL" else "generate_response"

workflow = StateGraph(AgentState)

# Add nodes
workflow.add_node("think", think)
workflow.add_node("execute_tool", execute_tool)
workflow.add_node("generate_response", generate_response)

# Add edges: a conditional edge after think, plain edges elsewhere
workflow.add_edge(START, "think")
workflow.add_conditional_edges(
    "think", route_after_think,
    {"execute_tool": "execute_tool", "generate_response": "generate_response"})
workflow.add_edge("execute_tool", "generate_response")
workflow.add_edge("generate_response", END)

app = workflow.compile()
result = app.invoke({"current_input": "What is 25 * 4 + 10?"})
print(result["final_answer"])
\end{lstlisting}

\end{frame}
